\section{Related Work}
\label{sec:related}

\paragraph{Lightweight image super-resolution.}
Lightweight SR has long focused on the tension between reconstruction quality and deployment cost. Earlier compact CNN models improve efficiency through post-upsampling, residual distillation, and lightweight feature fusion, while transformer-based variants inherit long-range interaction from self-attention~\cite{vaswani2017attention}. In practice, however, most lightweight transformer families still confine interaction to local windows or other geometry-defined neighborhoods in order to control quadratic complexity. This choice is efficient, but it also means that the interaction region is determined by spatial layout rather than visual similarity. For SR, this is a meaningful limitation because repeated structures and long contours frequently occur across distant but semantically related regions.

\paragraph{Content-aware token grouping.}
A promising direction is to organize tokens according to feature similarity before attention is applied. \baseline{}~\cite{liu2025catanet} is a representative example: it uses shared token centers to aggregate long-range content-similar tokens efficiently and then performs grouped interaction. Compared with content-agnostic windows, this strategy better exposes non-local redundancy in repeated structures and texture patterns. However, grouping alone is not enough. Once tokens have been gathered, the interaction stage must still decide how densely they should communicate, how subgroup impurity should be handled, and whether all channels should be modeled at the same restoration granularity. Our work targets exactly these remaining questions.

\paragraph{Focused and sparse attention.}
\pft{}~\cite{long2025pft} revisits the attention computation itself and argues that many token similarities are irrelevant to the current query. By linking attention maps across layers, progressive focused attention suppresses unimportant connections and reduces unnecessary computation. This perspective is especially attractive for SR, where repetitive textures and thin structures are usually supported by a few high-value correspondences rather than uniformly dense interactions. However, \pft{} is not built around content-aware grouping, and therefore it does not explicitly address the routing uncertainty, subgroup contamination, or scale allocation issues that emerge in grouped token interaction.

\paragraph{Our position.}
The proposed \method{} sits at the intersection of the two directions above. We retain the efficiency advantage of content-aware grouping, but we replace the single-scale dense interaction pattern with a multi-level focused sparse formulation. In our view, grouped token interaction should answer three questions jointly: which tokens should be routed together, which edges inside the subgroup graph should be kept, and how those edges should be distributed across low-, mid-, and high-frequency restoration branches. This combination yields a more complete answer to the lightweight SR problem than either content-aware grouping or focused attention alone.
